\section{Results}
\label{sec:results}
% RULE: every number in this section is read from a generated table or figure (tools/final_stats.py, tools/paper_extras.py,
% tools/bdcxr_stats.py -> paper/tables/, tools/fig_*.py and tools/paper_extras.py -> paper/figures/). Nothing is typed by hand without a source
% in results/final/stats_held.json or results/bdcxr/stats.json. Failed hypotheses are stated as failed with the numbers.
All numbers in this section are read from generated tables and figures. Of the pre-specified hypotheses, H1$'$ (dose-response) passes;
H2$'$ passes for bone-region PSNR against the reference model and against the no-twin control and fails for SSIM, as the analysis plan
anticipated; H3$'$ passes for both contrasts on VinDr-RibCXR; H4$'$a is mixed (ours above the reference model on JSRT and below it on Node21,
above the no-twin control on both); H4$'$b shows no downstream benefit of the soft-tissue image used alone; H5$'$ holds on synthetic data with
exact truth and fails under the shift to real radiographs. The v1 hypothesis H1 of the first analysis plan passes. All analysis changes made
after the freeze are listed in the deviations log of the repository.

\subsection{Fidelity on held-out paired radiographs (\texorpdfstring{H2$'$}{H2'})}
On the 51 held-out SZCH-X-Rays pairs (Table~\ref{tab:szch_heldout}) the bone-region PSNR$_{cal}$ of ours is 21.46\,dB and that of the
reference model 21.51\,dB. The paired difference is $-0.06$\,dB with a one-sided 95\,\% lower bound of $-0.15$\,dB, inside the
non-inferiority margin of $-0.5$\,dB (Table~\ref{tab:h2}). Against the no-twin control the difference is $+0.01$\,dB (lower bound
$-0.02$\,dB, margin $-0.25$\,dB): the twin regulariser costs no measurable bone-region fidelity. SSIM$_{cal}$ fails its margin: ours is
lower by 0.021 (lower bound $-0.023$, margin $-0.01$). The gap sits in fine texture. With a band limit of $\sigma=4$ the SSIM of ours and of
the reference model are 0.946 and 0.948, and the high-frequency energy of the non-bone region, relative to the dual-energy soft-tissue image,
is 1.09 for ours and 0.99 for the reference model; LPIPS is 0.234 and 0.166. BS-LDM, which was trained on dual-energy pairs from the same
hospital, reaches 28.57\,dB; it is shown for orientation and is not part of a pre-specified comparison. Figure~\ref{fig:pareto} places
fidelity against lesion absorption, and Figure~\ref{fig:qualitative} shows three held-out cases.
\begin{table}[htbp]
\centering\small
\caption{Held-out SZCH-X-Rays pairs: bone-region and whole-image metrics for ours and the comparators, calibrated (primary) and raw. HF
ratio: high-frequency energy of the non-bone region relative to the dual-energy soft-tissue image (1 = equal).}
\label{tab:szch_heldout}
\fittable{tables/szch_heldout.tex}
\end{table}
\begin{table}[htbp]
\centering\small
\caption{H2$'$ non-inferiority on the 51 held-out SZCH-X-Rays pairs (paired differences, image bootstrap, one-sided 95\,\% lower bound).}
\label{tab:h2}
\fittable{tables/h2_noninferiority.tex}
\end{table}
\begin{figure}[htbp]
\centering
\safefig{figures/pareto.pdf}
\caption{Lesion absorption on VinDr-RibCXR against bone-region PSNR$_{cal}$ on the 51 held-out SZCH-X-Rays pairs. Line: the SwinV2-T family at
twin weights $w=0\ldots64$ (single seeds light, seed means dark). Markers: ours and the no-twin control (ConvNeXt-V2-T, $K=3$), the reference
model, the v1 model and the identity. Lower absorption at equal fidelity is better.}
\label{fig:pareto}
\end{figure}
\begin{figure}[htbp]
\centering
\safefig{figures/qualitative_szch.png}
\caption{Held-out SZCH-X-Rays cases (CC BY-NC 4.0, \citealp{sun2026bsldm}): radiograph, dual-energy soft-tissue reference, the soft-tissue
estimates of ours and of the reference model, and the bone estimate of ours. The three cases are not hand-picked: they are the held-out images at
the 25th, 50th and 75th percentile of the bone-region PSNR$_{cal}$ of ours.}
\label{fig:qualitative}
\end{figure}

\subsection{Dose-response and ablation of the twin regulariser (\texorpdfstring{H1$'$}{H1'})}
\label{sec:res_twin}
In the SwinV2-T family (3, 2, 2, 4, 2 and 2 seeds at $w=0,4,8,16,32,64$; Table~\ref{tab:dose_response}, Figure~\ref{fig:dose_response}) the
mean absorption on VinDr-RibCXR falls from 0.183 at $w=0$ to 0.043 at $w=64$. The slope of $A$ on $\log_2(1+w)$ is $-0.0233$
$[-0.0249,-0.0217]$; its upper bound is below zero, so H1$'$ passes. On the Qure-projection test split with exact truth the slope is
$-0.0174$ $[-0.0190,-0.0160]$, with $A$ levelling off above $w=16$ (0.118, 0.122 and 0.112 at $w=16,32,64$). Bone-region PSNR$_{cal}$ on the
held-out pairs stays between 21.22 and 21.28\,dB up to $w=16$ and drops to 20.92 and 21.01\,dB at $w=32$ and 64, which is why $w=16$ is the
working point (Figure~\ref{fig:pareto}).

The ablation (Table~\ref{tab:ablation_twin}) is development evidence: it uses the audit on all 75 SZCH images and the 24 development pairs.
At $w=16$ the full regulariser gives $A=0.067$ against 0.190 without it. Three variants lose the effect: a stop-gradient on the branch without
the pasted lesion (0.185), a consistency term over the whole image instead of the lesion support (0.197), and lesion insertion without the twin
term (CLIA, 0.191). Dilating the support by 31\,px keeps part of it (0.137); up to four lesions per twin instead of up to two keep all of it (0.074).
The effect therefore needs gradients through both branches and a consistency term confined to where the lesion was pasted; showing the model
lesions during training is not enough.

The v1 hypothesis H1 of the first plan was tested on synthetic LIDC test CTs with exact truth on 27 October 2026, three seeds per arm
(Supplementary Table~S9), and passes: CLIA absorbs less than the recipe arm ($-0.022$ $[-0.032,-0.013]$, Holm $p<0.001$) with non-inferior bone
MAE. The lesion-weighted loss also lowers absorption ($-0.019$) and is non-inferior on the test CTs, although it missed the margin on the
validation CTs that decided the selection; combining both lowers absorption most ($-0.066$) but misses the margin. The 300k-sample finals of
CLIA and recipe replicate the contrast ($-0.024$ $[-0.034,-0.014]$). On real radiographs the v1 models remain far above the v2 models
($A=0.289$ and 0.308 on VinDr-RibCXR).
\begin{table}[htbp]
\centering\small
\caption{Dose-response of the lesion absorption on the twin weight (SwinV2-T family): seeds per weight, mean $A$ on VinDr-RibCXR and on the
Qure-projection test split, and the pre-specified slope with its image-cluster bootstrap interval. Each weight averages all of its seeds
(four at $w=16$, seeds 0--3); the encoder comparisons use seeds 0--2 only.}
\label{tab:dose_response}
\fittable{tables/dose_response.tex}
\end{table}
\begin{figure}[htbp]
\centering
\safefig{figures/dose_response.pdf}
\caption{Dose-response of the lesion absorption $A$ on VinDr-RibCXR on the twin weight $w$ (SwinV2-T), one marker per seed, with the
pre-specified slope on $\log_2(1+w)$.}
\label{fig:dose_response}
\end{figure}
\begin{table}[htbp]
\centering\small
\caption{Ablation of the twin regulariser at $w=16$ (SwinV2-T, 300k samples, mean over seeds): one factor changed per row. Development
evidence: audit on all 75 SZCH images, bone PSNR$_{cal}$ on the 24 development pairs.}
\label{tab:ablation_twin}
\fittable{tables/ablation_twin.tex}
\end{table}

\subsection{Confirmatory audit on VinDr-RibCXR (\texorpdfstring{H3$'$}{H3'})}
On 1,470 sites in 245 VinDr-RibCXR images (Tables~\ref{tab:audit_ribcxr} and~\ref{tab:audit_contrasts}) ours absorbs $A=0.051$
$[0.045,0.057]$ of an inserted lesion, the reference model 0.149 $[0.138,0.160]$ and the no-twin control 0.178 $[0.167,0.188]$. The paired
differences are $-0.098$ $[-0.108,-0.089]$ against the reference model and $-0.127$ $[-0.134,-0.119]$ against the no-twin control
(Holm-adjusted $p<0.001$ for both), so H3$'$ passes. The advantage does not depend on rib overlap here: the interaction between sites with at
least 15\,\% overlap with an annotated rib and the other sites is $0.003$ $[-0.017,0.027]$ against the reference model. On JSRT (secondary)
the pattern repeats ($A$ 0.048, 0.118 and 0.121; differences $-0.070$ and $-0.073$) with a larger advantage at bone-overlapping sites
(interaction $-0.060$ $[-0.074,-0.046]$); there the overlap was defined by the reference model's bone map, not by the consensus map the plan
names. The SZCH audit (development evidence) agrees ($A$ 0.054, 0.130 and 0.195). Figure~\ref{fig:audit_maps} shows one site.

Absorption must be read beside suppression strength, because doing nothing gives $A=0$ (Table~\ref{tab:rib_benchmark}). The reference model
removes more edge energy, both on the rib boundaries (48\,\% against 30\,\% for ours) and on the other lung pixels (22\,\% against 8\,\%);
its selectivity is 0.264 against 0.220 and the AUC separating rib from non-rib pixels is the same (0.671 and 0.672). Ours suppresses ribs
less strongly, but its absorption is a third of the reference model's while its rib-edge removal is about 60\,\% of it, and its bone-region
fidelity is non-inferior (H2$'$).
\begin{table}[htbp]
\centering\small
\caption{Counterfactual lesion audit on VinDr-RibCXR: mean absorption per method with image-cluster bootstrap intervals (lower is better).}
\label{tab:audit_ribcxr}
\fittable{tables/audit_ribcxr.tex}
\end{table}
\begin{table}[htbp]
\centering\small
\caption{H3$'$ paired site-level contrasts on the three audit sets, Holm over the two contrasts, and the interaction with bone overlap
($\geq$15\,\% vs.\ $<$15\,\%). SZCH rows: the 51 held-out images (306 sites), still development evidence because the audit on all 75 images
was used during development; the JSRT overlap comes from the reference model's bone map.}
\label{tab:audit_contrasts}
\fittable{tables/audit_contrasts.tex}
\end{table}
\begin{figure}[htbp]
\centering
\safefig{figures/audit_maps.pdf}
\caption{The audit on one held-out SZCH-X-Rays image (CC BY-NC 4.0): a 7\,mm nodule from a LIDC test CT is added to the radiograph; the change
of each bone estimate (red = lesion signal moved into the bone image) and the resulting soft-tissue image. Not hand-picked: eight placements
over the bone-overlapping lung field were drawn with a fixed seed and the one with the median $A(\text{ours})-A(\text{reference})$ is shown.}
\label{fig:audit_maps}
\end{figure}
\begin{table}[htbp]
\centering\small
\caption{Rib benchmark on VinDr-RibCXR (metrics only): gradient-energy reduction on rib-boundary pixels and on other lung pixels, selectivity
and AUC.}
\label{tab:rib_benchmark}
\fittable{tables/rib_benchmark.tex}
\end{table}

\subsection{Real nodules without lesion insertion (\texorpdfstring{H4$'$a}{H4'a})}
Real nodules give a mixed answer (Table~\ref{tab:nodule_cnr}). On JSRT (154 nodules) the contrast-to-noise ratio rises by 0.041 with ours
and by 0.018 with the reference model; the difference is $+0.023$ $[0.008,0.041]$, and $+0.032$ $[0.014,0.052]$ for the 123 nodules with at
least 15\,\% bone overlap. On Node21 (1,474 boxes) the reference model raises contrast more: $+0.103$ against $+0.078$, a difference of
$-0.025$ $[-0.033,-0.018]$, and $-0.017$ $[-0.026,-0.009]$ for the 1,109 bone-overlapping boxes. Against the no-twin control ours raises
real-nodule contrast in all four comparisons ($+0.007$ to $+0.027$, all intervals above zero; Holm within data set $p<0.001$).
\begin{table}[htbp]
\centering\small
\caption{Contrast-to-noise change of real nodules on JSRT and Node21 (soft-tissue image minus original), all lesions and lesions with at least
15\,\% bone overlap; image-cluster bootstrap intervals.}
\label{tab:nodule_cnr}
\fittable{tables/nodule_cnr.tex}
\end{table}

\subsection{Downstream tasks (\texorpdfstring{H4$'$b}{H4'b})}
The soft-tissue image used alone does not help downstream models (Tables~\ref{tab:downstream_cls} and~\ref{tab:downstream_det},
Figure~\ref{fig:downstream_forest}). On NIH ChestX-ray14 the macro-AUROC of the three-seed ensemble is 0.818 with the original image and
0.817 with ours ($-0.001$ $[-0.004,0.003]$, patient-cluster bootstrap); per class, ours is lower for effusion, infiltration and pleural
thickening ($-0.004$, $-0.007$, $-0.013$) and higher for nodule ($+0.011$; Holm-adjusted DeLong $p<0.05$ for these four). Adding the soft image
as a second channel helps slightly ($+0.006$ $[0.003,0.008]$), as does the reference model's soft image ($+0.004$ $[0.002,0.007]$). In
detection the change in CPM with ours alone is $+0.001$ $[-0.023,0.023]$ on Node21, $+0.006$ $[-0.009,0.020]$ on TBX11K and $-0.009$
$[-0.016,-0.002]$ on VinBigData; the three-channel input (original, soft and bone) gains $+0.028$ $[0.009,0.047]$ on Node21 and $+0.008$
$[0.002,0.014]$ on VinBigData, and the reference model's soft image gains $+0.029$ $[0.016,0.043]$ on TBX11K. No difference in differences of
the recall at two false positives per image between bone-overlapping and other lesions excludes zero (Node21 $-0.006$ $[-0.071,0.061]$,
VinBigData $+0.009$ $[-0.034,0.053]$). On TBX11K it cannot be estimated, because every test box overlaps bone by at least 20\,\%.
\begin{table}[htbp]
\centering\small
\caption{Downstream classification on NIH ChestX-ray14: macro-AUROC of the three-seed ensemble, per-seed mean $\pm$ SD and the difference to
the original image (patient-cluster bootstrap).}
\label{tab:downstream_cls}
\fittable{tables/downstream_cls.tex}
\end{table}
\begin{table}[htbp]
\centering\small
\caption{Downstream detection on Node21, TBX11K and VinBigData: CPM and mAP@50 (mean $\pm$ SD over three seeds), differences to the original
image and the difference in differences (DiD) of recall at two false positives per image (image bootstrap). both: the original, soft-tissue
and bone images as three channels.}
\label{tab:downstream_det}
\fittable{tables/downstream_det.tex}
\end{table}
\begin{figure}[htbp]
\centering
\safefig{figures/downstream_forest.pdf}
\caption{Downstream effects of each input against the original image with 95\,\% intervals: macro-AUROC on NIH ChestX-ray14, CPM on the
detection sets.}
\label{fig:downstream_forest}
\end{figure}

\subsection{Reliability (\texorpdfstring{H5$'$}{H5'})}
On the Qure-projection test split with exact truth the split-conformal intervals of the three-model ensemble cover 0.900 $[0.897,0.903]$ of
the pixels at a nominal 0.90, and the ensemble spread ranks the absolute bone error well (Spearman 0.76; Table~\ref{tab:reliability},
Figure~\ref{fig:reliability}). On the 51 held-out SZCH-X-Rays pairs, measured against a pseudo-truth bone image from the dual-energy pair, the
same intervals cover only 0.139 $[0.134,0.145]$ and the spread ranks the error weakly (Spearman 0.11 $[0.07,0.15]$). The uncertainty is
calibrated in the training domain and does not transfer to real radiographs without recalibration.
\begin{table}[htbp]
\centering\small
\caption{Ensemble spread and split-conformal intervals ($K=3$, nominal coverage 0.90), with image-bootstrap intervals.}
\label{tab:reliability}
\fittable{tables/reliability.tex}
\end{table}
\begin{figure}[htbp]
\centering
\safefig{figures/reliability.pdf}
\caption{Left: per-image coverage of the conformal intervals on the Qure-projection test split and on the held-out SZCH-X-Rays pairs (bars:
means; dashed: nominal 0.90). Right: per-image Spearman correlation of ensemble spread and absolute error on SZCH-X-Rays (line: mean).}
\label{fig:reliability}
\end{figure}

An exploratory analysis, added after this result was known and therefore post hoc on already-seen pairs, recalibrated the intervals on the
24 SZCH-X-Rays development pairs and tested them on the 51 held-out pairs (Table~\ref{tab:bone_recal}, Supplementary Fig.~S10). Pooled
coverage was 0.886 with one global calibration and 0.895 with a separate calibration per stratum of predicted bone overlap and ensemble spread
(Mondrian), whose intervals were less than half as wide (mean width 0.215 against 0.447). A few real pairs thus restore nominal coverage on
average, while the coverage of single images still varies.
\begin{table}[htbp]
\centering\small
\caption{Exploratory recalibration on real pairs (post hoc): bone-image intervals calibrated on the 24 SZCH-X-Rays development pairs and tested
on the 51 held-out pairs at nominal 0.90. Mondrian: calibration per stratum of predicted bone overlap and ensemble spread.}
\label{tab:bone_recal}
\fittable{tables/v3_bone_intervals.tex}
\end{table}
\subsection{Encoders (exploratory)}
The twin effect replicates across three encoders trained under the same recipe (Table~\ref{tab:encoders}). On VinDr-RibCXR the twin lowers
$A$ by 0.127 for ConvNeXt-V2-T, 0.114 for SwinV2-T and 0.112 for RAD-DINO, and on JSRT by 0.073, 0.061 and 0.064. ConvNeXt-V2-T was chosen
as ``ours'' on development data before the freeze (Section~\ref{sec:experiments}); RAD-DINO reaches the same absorption ($A=0.053$ on
VinDr-RibCXR).
\begin{table}[htbp]
\centering\small
\caption{Twin regulariser per encoder (exploratory): absorption difference to the matched no-twin control and its interaction with bone
overlap; three seeds per model (seeds 0--2), SZCH on its 51 held-out images.}
\label{tab:encoders}
\fittable{tables/encoders.tex}
\end{table}

\subsection{External validation on paediatric radiographs (BDCXR-3257)}
\label{sec:bdcxr}
On 3,257 paediatric radiographs from a hospital in Dhaka (Table~\ref{tab:bdcxr_cls}), five-fold cross-validated pneumonia classification
(normal, bacterial, viral) reaches a pooled out-of-fold macro-AUROC of 0.814 on the original image and 0.816 on the soft-tissue image of
ours. The difference, $+0.002$ $[-0.005,0.010]$, meets the pre-specified non-inferiority margin of 0.01 and shows no superiority; the
reference model's soft image gives $+0.001$ $[-0.008,0.010]$ and CLAHE $-0.006$ $[-0.012,0.001]$, which misses the margin. In the lesion audit
on 300 of these images (Table~\ref{tab:bdcxr_audit}) ours absorbs $A=0.084$, the reference model 0.217 and the no-twin control 0.251
(differences $-0.134$ $[-0.144,-0.124]$ and $-0.167$ $[-0.175,-0.160]$), with a larger advantage at bone-overlapping sites. The inserted
lesions come from adult CTs, so this is a stress test rather than a clinical finding. The ensemble disagreement on these children is 2.7 times
that on the adult SZCH-X-Rays pairs (median lung-field spread 0.0076 against 0.0028) and is larger in younger children (Spearman with age
$-0.22$).
\begin{table}[htbp]
\centering\small
\caption{BDCXR-3257: pneumonia classification (normal, bacterial, viral), five folds over near-duplicate groups $\times$ three seeds;
pooled out-of-fold macro-AUROC, per-fold mean $\pm$ SD and the difference to the original image (group bootstrap).}
\label{tab:bdcxr_cls}
\fittable{tables/bdcxr_cls.tex}
\end{table}
\begin{table}[htbp]
\centering\small
\caption{BDCXR-3257: lesion audit on a fixed class-stratified subset of 300 images (1,800 sites).}
\label{tab:bdcxr_audit}
\fittable{tables/bdcxr_audit.tex}
\end{table}
